\chapter{排序}
\label{chap:sorting}

\setcounter{footnote}{0}

\begin{flushright}
  \small Michael Garland 特别贡献
\end{flushright}

排序算法把列表中的数据元素排列成特定次序。如果数据集处于适当次序，从中检索信息的
计算复杂度就能显著降低，因此排序是现代数据与信息服务的基础。例如，排序经常用于把
数据转换为规范形式，以便快速比较并协调不同数据列表；许多数据处理算法也能因数据具有
特定次序而提高效率。高效排序算法的重要性吸引了大量计算机科学研究。即使采用这些高效
算法，对大型数据列表排序仍然很耗时，可以从并行执行中受益。高效排序算法的并行化颇具
挑战，有时需要精巧的设计。

本章介绍三种并行排序算法的设计：奇偶排序、归并排序和基数排序。奇偶排序十分简单，
因此只作简要介绍；归并排序主要依赖第 13 章已经介绍的归并模式，也只作简述。本章大部分
篇幅用于基数排序及其优化；这些方法依赖第 12 章稳定过滤模式的一种更一般形式。最后还会
讨论双调排序、采样排序等其他流行的并行排序算法。

\section{背景}
\label{sec:sorting-background}

排序是计算机最早的应用之一。排序算法把列表元素排列成特定次序。应当采用什么次序，
取决于元素本身的性质；常见例子包括数值的数值顺序和文本字符串的字典序。更形式化地说，
任何排序算法都必须满足以下两个条件：

\begin{itemize}
  \item 输出按非递减或非递增次序排列。对非递减次序，按照目标顺序，每个元素都不小于
  前一个元素；对非递增次序，每个元素都不大于前一个元素。
  \item 输出是输入的一个排列。也就是说，算法在重新排列输入元素形成输出的同时，必须
  保留原输入中的所有元素。
\end{itemize}

最简单的情况下，列表元素可以按元素自身的值排序。例如，列表
$[5,2,7,1,3,2,8]$ 可以按非递减次序排序为 $[1,2,2,3,5,7,8]$。一种更复杂也更
常见的情况是，每个元素由键字段和值字段组成，列表按键字段排序。假设每个元素都是
二元组（年龄，以千美元计的收入），列表
$[(30,150),(32,80),(22,45),(29,80)]$ 以收入为键按非递增次序排序后得到
$[(30,150),(32,80),(29,80),(22,45)]$。如果列表元素本身只是数值，那么排序时
这些数值就充当键。

排序算法可以分为稳定算法与不稳定算法。当两个元素的键值相等时，\textbf{稳定排序
（stable sort）}会保留它们在输入中出现的先后次序。例如，以收入为键对上述列表按
非递增次序排序时，稳定排序算法必须保证 $(32,80)$ 出现在 $(29,80)$ 之前，因为前者
在原输入中先出现；不稳定排序则不提供这种保证。若要用多个键以级联方式对列表排序，
就需要稳定算法。例如，每个元素都有主键与次键时，可以先按次键排序，再按主键排序；
第二次稳定排序会保留第一次排序在相同主键元素之间建立的次序。

排序算法还可以分为基于比较与非基于比较的算法。对包含 $N$ 个元素的列表，基于比较的
排序算法必须完成一定数量的元素间比较，因此复杂度不可能优于 $O(N\log N)$。相比之下，
某些非比较排序算法能取得优于 $O(N\log N)$ 的复杂度，但它们未必适用于任意类型的键。
两类算法都可以并行化。本章将介绍两种基于比较的并行排序算法，即奇偶排序和归并排序，
以及一种非比较并行排序算法，即基数排序。

排序的重要性促使计算机科学研究界基于丰富的数据结构与算法策略，发展出了种类繁多的
排序算法。因此，计算机科学导论课程经常用排序算法说明许多核心算法概念，包括大 $O$
记号、分治算法、堆和二叉树等数据结构、随机算法、最好/最坏/平均情况分析、时间与空间
权衡，以及上界与下界。本章延续这一传统，用三种排序算法说明若干重要的并行化和性能
优化技术\cite{satish2009sorting}。

\section{并行奇偶排序}
\label{sec:parallel-odd-even-sort}

冒泡排序是最简单的基于比较的顺序排序算法之一。它的高层策略是比较列表中相邻的元素
对；若顺序错误，就交换二者。不断重复这一过程，直到整个列表有序。冒泡排序的每次迭代
都从头到尾顺序遍历元素列表并整理相邻元素对。不过，只要每次迭代覆盖所有相邻对，处理
各对的次序并非算法正确性的必要条件。并行排序算法可以采用相似的迭代策略，但在每次
迭代中并行而非顺序地整理各相邻对。

并行处理元素对的挑战在于，被同时处理的各对必须相互独立，才能避免竞争条件。例如，
若一个线程负责排序下标 0 和 1 处的元素，就不能再让另一个线程同时排序下标 1 和 2
处的元素。如果两个线程同时决定交换各自元素对，访问下标 1 处的元素时便会发生竞争。
为了最大化同时处理的独立元素对数，可以先并行排序首元素下标为偶数的所有元素对
（偶数对），再并行排序首元素下标为奇数的所有元素对（奇数对）。这种技术构成了
\textbf{奇偶换位排序（odd-even transposition sort）}的基础，简称奇偶排序。

奇偶排序是一种基于比较的并行排序算法，它在相邻元素的奇数对与偶数对之间交替迭代。
每次迭代比较所有奇数对或偶数对，若顺序错误就交换其中元素。图 \ref{fig:14-1} 给出
对八个元素执行奇偶排序的例子。第一次迭代启动四个线程，每个线程负责一个偶数对。
前两个线程发现其元素对已经有序，后两个线程发现顺序错误并交换元素。第二次迭代为每个
奇数对启动一个线程，同样只有部分线程执行交换。第三次迭代重新处理偶数对。算法持续
交替处理奇数对与偶数对，直到没有任何元素被交换，表明列表已有序。

\begin{figure}[htbp]
  \centering
  \small
  \setlength{\tabcolsep}{5pt}
  \begin{tabular}{c|*{8}{c}}
    初始&4&7&2&3&8&5&9&6\\
    整理偶数对&4&7&2&3&5&8&6&9\\
    整理奇数对&4&2&7&3&5&6&8&9\\
    整理偶数对&2&4&3&7&5&6&8&9\\
    整理奇数对&2&3&4&5&7&6&8&9\\
    整理偶数对&2&3&4&5&6&7&8&9\\
    整理奇数对（无变化）&2&3&4&5&6&7&8&9
  \end{tabular}
  \caption{并行奇偶排序}
  \label{fig:14-1}
\end{figure}

图 \ref{fig:14-2} 给出奇偶排序内核。列表中有多少元素，每次迭代就以其一半的线程数
调用一次内核。参数 \code{isOddStep}（第 2 行）指明当前内核处理奇数对还是偶数对。
内核首先确定线程负责的元素对中第一个元素的下标（第 3--4 行）。对偶数对
（\code{isOddStep=0}），该下标就是线程全局下标的两倍，因此线程 $t$ 处理下标
$(2t,2t+1)$；对奇数对（\code{isOddStep=1}），再加 1，于是线程 $t$ 处理下标
$(2t+1,2t+2)$。通过边界检查后（第 5 行），每个线程比较自己负责的两个元素（第 6
行），若顺序错误就交换（第 7--9 行）。若交换成功，线程还会设置全局标志
\code{hasChanged}，表示列表次序已经变化，需要继续执行一次迭代（第 10 行）。

\begin{figure}[htbp]
  \centering
  \begin{minipage}{0.96\textwidth}
  \begin{lstlisting}[language=C++,basicstyle=\ttfamily\scriptsize]
__global__ void sort_kernel(unsigned int* data, unsigned int* hasChanged,
                            unsigned int N, unsigned int isOddStep) {
  unsigned int i = 2*(blockIdx.x*blockDim.x + threadIdx.x)
                   + (isOddStep ? 1 : 0);
  if(i < N - 1) {
    if(data[i] > data[i + 1]) {
      unsigned int tmp = data[i];
      data[i] = data[i + 1];
      data[i + 1] = tmp;
      *hasChanged = 1;
    }
  }
}
  \end{lstlisting}
  \end{minipage}
  \caption{并行奇偶排序内核}
  \label{fig:14-2}
\end{figure}

多个线程同时向 \code{hasChanged} 写入 1 时，从严格意义上说存在竞争条件：多个线程
无序访问同一内存位置，且至少一次访问为写入。实践中，这类竞争通常无害，因为所有线程
都写入同一个值，无论多少线程执行写入，最终结果都相同。一个操作无论应用多少次效果都
相同的性质称为\textbf{幂等性（idempotence）}。然而原则上，即使操作幂等，C++ 内存
模型中的竞争条件仍可能导致未定义行为。因此，这段代码在实践中或许能工作，却违反
C++ 内存模型，不能保证正确。保守的程序员应当使用原子操作向 \code{hasChanged}
写入 1。

奇偶排序易于实现，但时间复杂度和工作复杂度都很高。所需迭代次数为 $O(N)$，其中 $N$
为输入列表大小；因为最坏情况下最大元素位于列表开头，需要 $N$ 次迭代才能抵达末尾。
若执行资源充足，使每次迭代都能完全并行执行，则时间复杂度为 $O(N)$。每次迭代执行
$N/2$ 次，即 $O(N)$ 次比较，因此奇偶排序的工作复杂度为 $O(N^2)$。

\section{并行归并排序}
\label{sec:parallel-merge-sort}

奇偶排序的时间与工作复杂度很高，因此需要复杂度更低、更高效的基于比较的并行排序算法。
归并排序就是一种适合并行化的比较排序算法。它先把输入列表划分为若干区段，分别用归并
排序或其他算法排好每个区段，再对有序区段执行有序归并。

图 \ref{fig:14-3} 展示了并行化归并排序的一种方法。开始时，把输入列表划分为许多
区段，每个区段用某种排序算法独立排序。随后，把每对区段归并成一个区段。不断重复这一
过程，直到所有键都归入同一区段。在每一阶段，既可以并行执行不同归并操作，也可以从每个
归并操作内部提取并行性。第 13 章已经介绍了如何并行化归并操作。

\begin{figure}[htbp]
  \centering
  \small
  \begin{tabular}{c}
    \fbox{块 0}\ \fbox{块 1}\ \fbox{块 2}\ \fbox{块 3}\
    \ \fbox{块 4}\ \fbox{块 5}\ \fbox{块 6}\ \fbox{块 7}\\[0.3em]
    $\downarrow$ 各块独立排序 $\downarrow$\\[0.3em]
    \fbox{有序段 0}\ \fbox{有序段 1}\ \fbox{有序段 2}\ \fbox{有序段 3}\
    \ \fbox{有序段 4}\ \fbox{有序段 5}\ \fbox{有序段 6}\ \fbox{有序段 7}\\[0.3em]
    $\searrow\!\swarrow\quad\searrow\!\swarrow\quad
     \searrow\!\swarrow\quad\searrow\!\swarrow$\\[-0.1em]
    \fbox{归并段 0--1}\quad\fbox{归并段 2--3}\quad
    \fbox{归并段 4--5}\quad\fbox{归并段 6--7}\\[0.3em]
    $\searrow\!\swarrow\qquad\qquad\searrow\!\swarrow$\\[-0.1em]
    \fbox{归并段 0--3}\qquad\qquad\fbox{归并段 4--7}\\[0.3em]
    $\searrow\qquad\swarrow$\\[-0.1em]
    \fbox{完整有序列表}
  \end{tabular}
  \caption{并行化归并排序的一种方式}
  \label{fig:14-3}
\end{figure}

随着图 \ref{fig:14-3} 的并行归并排序逐阶段推进，从归并操作之间与归并操作内部可
提取的并行性之间存在权衡。较早阶段有更多独立归并操作可以并行执行；较晚阶段的独立
归并操作较少，但每个操作归并更多键，因此内部暴露出更多并行性。例如，第一个归并阶段
包含四个独立归并操作，一个由八个线程块构成的网格可以为每个归并操作分配两个线程块。
下一阶段只有两个归并操作，但每个操作处理两倍的键，因此可以为每个操作分配四个线程块。
这一并行归并排序的实现留作练习。

图 \ref{fig:14-3} 的并行归并排序需要 $O(\log N)$ 次迭代，因为每次迭代都把待归并
区段数减半。若执行资源充足，使每次迭代中的归并操作都能完全并行执行，则按照第 13 章
的分析，每次迭代的时间复杂度为 $O(\log N)$，工作复杂度为 $O(N\log N)$。因此，
这里的并行归并排序总时间复杂度和总工作复杂度分别为 $O(\log^2 N)$ 和
$O(N\log^2 N)$。请注意，$O(N\log^2 N)$ 略高于顺序归并排序的 $O(N\log N)$；
原因是第 13 章所示的并行归并需要执行协同秩（co-rank）函数，工作复杂度高于顺序归并。
Cole 的并行归并排序算法\cite{cole1988merge}克服了这一限制，分别取得
$O(\log N)$ 的时间复杂度和 $O(N\log N)$ 的工作复杂度。

\section{基数排序}
\label{sec:radix-sort}

要取得优于 $O(N\log N)$ 的排序复杂度，必须使用非比较排序算法。基数排序是一种非常
适合并行化的非比较排序算法。它根据基数值，即位值计数制中的底数，把待排序键分配到
不同\textbf{桶（bucket）}中。如果键由多个数位组成，就对每一位重复分配，直到处理
全部数位。每次迭代都是稳定的，会在每个桶内保留上次迭代形成的键次序。处理二进制表示
的键时，选择 2 的幂作为基数很方便，因为这样易于遍历并提取数位；每次迭代本质上处理
键中一个固定宽度的位片。本章先使用基数 2，即 1 位基数，之后再扩展到更大基数。

\begin{figure}[htbp]
  \centering
  \tiny
  \setlength{\tabcolsep}{1.5pt}
  \begin{tabular}{c|*{16}{c}}
    输入&1100&0011&0110&1001&1111&1000&0101&1010&1001&0110&1011&1101&0100&1010&0111&0000\\
    第 1 轮（位 0）&1100&0110&1000&1010&0110&0100&1010&0000&0011&1001&1111&0101&1001&1011&1101&0111\\
    第 2 轮（位 1）&1100&1000&0100&0000&1001&0101&1001&1101&0110&1010&0110&1010&0011&1111&1011&0111\\
    第 3 轮（位 2）&1000&0000&1001&1001&1010&1010&0011&1011&1100&0100&0101&1101&0110&0110&1111&0111\\
    第 4 轮（位 3）&0000&0011&0100&0101&0110&0110&0111&1000&1001&1001&1010&1010&1011&1100&1101&1111
  \end{tabular}
  \caption{1 位基数的基数排序示例}
  \label{fig:14-4}
\end{figure}

图 \ref{fig:14-4} 展示了如何用 1 位基数对 4 位整数列表进行基数排序。键长为 4 位，
每次迭代处理 1 位，所以总共需要四次迭代。第一次迭代考察最低有效位。最低位为 0 的
所有键放到输出列表左侧，构成 0 桶；最低位为 1 的键放到右侧，构成 1 桶。输出列表的
每个桶内都保留输入列表中的键次序。换句话说，进入同一桶的键在输出中必须保持它们在
输入中的相对顺序。讨论下一次迭代时，就会看到为什么稳定性要求十分重要。

第二次迭代把第一次迭代的输出作为新输入，并考察每个键的次低位。与第一次相同，键被
分入次低位为 0 和为 1 的两个桶。由于保留了前一次迭代的次序，第二次迭代的输出已经
按低 2 位有序：低 2 位为 00 的键在最前，然后依次为 01、10 和 11。

第三次迭代对键的第 3 位重复这一过程。由于仍保留前几次迭代的次序，输出列表按低 3 位
有序。最后一次迭代考察第 4 位，即最高有效位。迭代结束后，最终输出列表按全部 4 位
完全有序。

基数排序一次迭代执行的操作，应当使读者想到第 12 章的稳定过滤模式。稳定过滤把满足
某个条件的输入键搬到输出列表开头，同时保持它们的次序。基数排序迭代执行一种更一般的
操作：满足条件，即当前位为 0 的输入键被搬到输出列表开头；不满足条件的键被搬到列表
末尾，同时两个子列表内部都保持键的次序。这一更一般的模式有时称为\textbf{稳定划分
（stable partition）}，因为它把输入列表分成多个分区。下文将看到，基数排序迭代的
并行化与优化技术同稳定过滤非常相似。

\section{并行基数排序}
\label{sec:parallel-radix-sort}

基数排序的每次迭代都依赖上一次迭代的完整结果，因此各次迭代必须依次执行。并行化机会
存在于每次迭代内部。本章余下部分集中讨论单次迭代的并行化与优化，并默认主机代码为
每次迭代调用一次这样的内核。我们将介绍该基数排序内核的多项优化。感兴趣的读者还可
参考 OneSweep\cite{adinets2022onesweep}，它是目前先进的 GPU 基数排序实现。

在 GPU 上并行化一次基数排序迭代，一种直接方法是让每个线程负责输入列表中的一个键。
线程先确定自己的键在输出列表中的位置，再把键存到该位置。图 \ref{fig:14-5} 把这种
方法应用于图 \ref{fig:14-4} 的第一次迭代。概念图以四个线程块处理 16 个键，每块
四个线程，每个线程负责输入列表中对应的键。实际情况下，每块最多可以有 1024 个线程，
输入也远大于本例，因而会产生更多线程块；这里使用小规模设置只是为了便于说明。

\begin{figure}[htbp]
  \centering
  \scriptsize
  \begin{tabular}{c}
    \fbox{块 0：1100 0011 0110 1001}\quad
    \fbox{块 1：1111 1000 0101 1010}\\
    \fbox{块 2：1001 0110 1011 1101}\quad
    \fbox{块 3：0100 1010 0111 0000}\\[0.4em]
    $\downarrow$ 每个线程根据最低位求目标下标 $\downarrow$\\[0.4em]
    \fbox{0 桶：1100 0110 1000 1010 0110 0100 1010 0000}\\
    \fbox{1 桶：0011 1001 1111 0101 1001 1011 1101 0111}
  \end{tabular}
  \caption{为每个输入键分配一个线程，并行化一次基数排序迭代}
  \label{fig:14-5}
\end{figure}

每个线程取得一个输入键后，剩余挑战是确定它在输出列表中的目标下标。求目标下标的方法
取决于该键进入 0 桶还是 1 桶。对进入 0 桶的键，有

\begin{align*}
  \text{0 键的目标下标}
  &=\text{此前 0 键数}\\
  &=\text{此前键总数}-\text{此前 1 键数}\\
  &=\text{键下标}-\text{此前 1 键数}.
\end{align*}

进入 0 桶的键的目标下标，等于它之前同样进入 0 桶的键数。每个键非入 0 桶即入 1 桶，
所以此前 0 键数等于此前键总数减去此前 1 键数。此前键总数就是该键在输入列表中的下标，
可以直接取得。因此，唯一不平凡的步骤是统计它之前进入 1 桶的键数；稍后将用排他扫描
完成这一操作。

对进入 1 桶的键，有

\begin{align*}
  \text{1 键的目标下标}
  &=\text{0 键总数}+\text{此前 1 键数}\\
  &=(\text{键总数}-\text{1 键总数})+\text{此前 1 键数}\\
  &=\text{输入大小}-\text{1 键总数}+\text{此前 1 键数}.
\end{align*}

输出数组中，所有进入 0 桶的键都必须位于进入 1 桶的键之前。因此，1 键的目标下标等于
0 键总数加此前 1 键数。0 键总数又等于输入键总数减去 1 键总数，而输入键总数就是可
直接取得的输入大小。于是，这里同样需要统计此前 1 键数；排他扫描还能顺便给出 1 键
总数。

\begin{figure}[htbp]
  \centering
  \tiny
  \setlength{\tabcolsep}{2.1pt}
  \begin{tabular}{c|*{16}{c}}
    键下标&0&1&2&3&4&5&6&7&8&9&10&11&12&13&14&15\\
    当前位&0&1&0&1&1&0&1&0&1&0&1&1&0&0&1&0\\
    此前 1 键数&0&0&1&1&2&3&3&4&4&5&5&6&7&7&7&8\\
    目标下标&0&8&1&9&10&2&11&3&12&4&13&14&5&6&15&7
  \end{tabular}
  \caption{求每个输入键的目标位置}
  \label{fig:14-6}
\end{figure}

图 \ref{fig:14-6} 展示了本例中每个线程为确定目标下标而执行的操作，相应内核见图
\ref{fig:14-7}。每个线程先确定自己负责的键下标（第 3 行），进行边界检查（第 5
行），再从输入列表加载该键（第 6 行）。随后，线程从键中提取本轮考察的位，确定它是
0 还是 1（第 7 行）。迭代号 \code{iter} 给出目标位的位置。把键右移该位数，就把
目标位移到最右端；再让移位后的键与 1 执行按位与，就把最低位以外的所有位清零，
\code{bit} 最终得到目标位的值。

\begin{figure}[htbp]
  \centering
  \begin{minipage}{0.98\textwidth}
  \begin{lstlisting}[language=C++,basicstyle=\ttfamily\scriptsize]
__global__ void radix_sort_iter(unsigned int* input, unsigned int* output,
                                unsigned int* bits, unsigned int N,
                                unsigned int iter) {
  unsigned int i = blockIdx.x*blockDim.x + threadIdx.x;
  unsigned int key, bit;
  if(i < N) {
    key = input[i];
    bit = (key >> iter) & 1;
    bits[i] = bit;
  }
  gridExclusiveScan(bits, N);
  if(i < N) {
    unsigned int numOnesBefore = bits[i];
    unsigned int numOnesTotal = bits[N];
    unsigned int dst = (bit == 0) ? (i - numOnesBefore)
                                  : (N - numOnesTotal + numOnesBefore);
    output[dst] = key;
  }
}
  \end{lstlisting}
  \end{minipage}
  \caption{基数排序单次迭代内核代码}
  \label{fig:14-7}
\end{figure}

每个线程提取目标位后，把它写入内存（第 8 行），各线程再协作对这些位执行网格范围的
排他扫描（第 10 行）。第 11 章已经介绍了网格级并行扫描。排他扫描结果数组的每个位置
保存该位置之前所有位的和。位值只能是 0 或 1，因此这个和就等于此前 1 的数量，也就是
此前进入 1 桶的键数。每个线程读取数组，取得自身位置之前的 1 键数（第 12 行）和输入
列表中的 1 键总数（第 13 行），再用前面推导的表达式确定目标位置（第 14--15 行）。
最后，线程把自己的键写到输出列表中的相应位置（第 16 行）。

\noindent\textbf{译者注：}图 \ref{fig:14-7} 读取 \code{bits[N]} 作为 1 键总数，
因此调用方需要为 \code{bits} 分配至少 $N+1$ 个元素，并保证
\code{gridExclusiveScan(bits,N)} 按底本约定把总和写到该位置。该接口是本章用来
表达算法结构的抽象调用，具体实现应与第 11 章扫描内核的接口约定保持一致。

\section{针对内存合并进行优化}
\label{sec:radix-memory-coalescing}

上述方法能有效并行化基数排序迭代，但一个主要低效来源是：对输出列表的写入访问模式
难以充分合并。观察图 \ref{fig:14-5} 中各线程如何写入输出列表。第一个线程块中，
第一个线程写 0 桶，第二个写 1 桶，第三个写 0 桶，第四个写 1 桶。因此，连续下标的
线程不一定写入连续内存位置，内存合并效果较差，每个 warp 需要发出多个内存请求。

对只有两个桶的 1 位基数，同一 warp 中的线程只会写入这两个桶之一。因此，一个 warp
写入的不同缓存行通常为两条，或因访问对齐而略多。但第 \ref{sec:radix-choice} 节将
看到，更大的基数会产生更多桶，同一 warp 可能写入更多不同缓存行。因此，未合并的
内存访问会严重妨碍基数排序迭代内核的性能。

回忆第 6 章，改进内核内存合并有多种办法：（1）重新排列线程；（2）重新排列线程访问
的数据；（3）在共享内存上执行无法合并的访问，再以合并方式在共享内存和全局内存之间
传输数据。本章采用第三种方法。各线程不再以难以合并的方式直接把键写到全局内存桶，
而是在共享内存中为每个线程块维护局部桶。也就是说，不再像图 \ref{fig:14-7} 那样
直接执行全局排序。块内线程先完成块级局部排序，在共享内存中把进入 0 桶和 1 桶的键
分开，再把各桶从共享内存以合并方式写入全局内存。

\begin{figure}[htbp]
  \centering
  \scriptsize
  \begin{tabular}{c}
    \fbox{块 0：1100 0011 0110 1001 $\rightarrow$ 1100 0110 $\mid$ 0011 1001}\\
    \fbox{块 1：1111 1000 0101 1010 $\rightarrow$ 1000 1010 $\mid$ 1111 0101}\\
    \fbox{块 2：1001 0110 1011 1101 $\rightarrow$ 0110 $\mid$ 1001 1011 1101}\\
    \fbox{块 3：0100 1010 0111 0000 $\rightarrow$ 0100 1010 0000 $\mid$ 0111}\\[0.35em]
    $\downarrow$ 共享内存中先按块局部排序，再连续写出局部桶 $\downarrow$\\[0.35em]
    \fbox{全局 0 桶：1100 0110 1000 1010 0110 0100 1010 0000}\\
    \fbox{全局 1 桶：0011 1001 1111 0101 1001 1011 1101 0111}
  \end{tabular}
  \caption{先在共享内存中局部排序，再写入全局内存，以改进内存合并}
  \label{fig:14-8}
\end{figure}

图 \ref{fig:14-8} 展示了如何改进图 \ref{fig:14-5} 示例中的内存合并。每个线程块
先对自己拥有的键执行局部基数排序，并把输出列表存入共享内存。局部排序方式与先前的
全局排序相同，只需每个线程块执行局部排他扫描，而不需要全局扫描。局部排序后，各块以
更易合并的方式把局部桶写入全局桶。例如，第一个线程块写出 0 桶时，前两个线程写入
全局内存相邻位置；写出 1 桶时，后两个线程也写入相邻位置。因此，大部分全局内存写入
都能合并。

该优化的主要挑战，是每个线程块必须确定自己的每个局部桶在对应全局桶中的起始位置。
局部桶的起始位置取决于其他线程块中局部桶的大小。具体来说，某块的局部 0 桶位于所有
先前线程块局部 0 桶之后；该块的局部 1 桶则位于所有线程块的局部 0 桶以及先前线程块
的局部 1 桶之后。对各线程块的局部桶大小执行排他扫描，就能得到这些位置。

\begin{figure}[htbp]
  \centering
  \small
  \setlength{\tabcolsep}{8pt}
  \begin{tabular}{c|cccc}
    线程块&0&1&2&3\\\hline
    局部 0 桶大小&2&2&1&3\\
    局部 1 桶大小&2&2&3&1\\\hline
    \multicolumn{5}{c}{$\Downarrow$ 对按行展开的桶大小表执行排他扫描 $\Downarrow$}\\\hline
    局部 0 桶目标起点&0&2&4&5\\
    局部 1 桶目标起点&8&10&12&15
  \end{tabular}
  \caption{确定每个线程块的局部桶在全局输出中的目标位置}
  \label{fig:14-9}
\end{figure}

图 \ref{fig:14-9} 展示了如何用排他扫描求每个线程块的局部桶位置。局部基数排序完成
后，各线程块确定每个局部桶中的键数，再把这些值存入图示表格。表格采用行主序：所有
线程块的局部 0 桶大小连续排列在前，随后是所有局部 1 桶大小。表格建立后，对线性化
表格执行排他扫描；结果表格就是每个线程块局部桶的起始位置。

线程块知道局部桶在全局内存中的起点后，块内线程便可把键从局部桶写到全局桶。为此，
每个线程需要掌握局部 0 桶与局部 1 桶的键数。写入阶段中，块内线程根据自己的线程下标
写其中一个桶。例如图 \ref{fig:14-9} 的块 2 中，线程 0 写 0 桶中唯一的键，线程
1--3 写 1 桶中的三个键；块 3 中则由线程 0--2 写 0 桶的三个键，线程 3 写 1 桶
的一个键。每个线程需要判断自己负责写局部 0 桶还是 1 桶。每块记录两个局部桶中的键数，
让线程能据此确定自己的 \code{threadIdx} 落在哪个范围，并参与相应桶的写入。该优化
的具体实现留作练习。

\section{基数值的选择}
\label{sec:radix-choice}

目前一直以 1 位基数说明如何并行化基数排序。对示例中的 4 位键，需要四次迭代才能
完全排序；一般而言，$N$ 位键需要 $N$ 次迭代。使用更大的基数值可以减少迭代次数。

图 \ref{fig:14-10} 展示了使用 2 位基数的基数排序。每次迭代用 2 位把键分入各桶，
所以只需两次迭代就能把 4 位键完全排好。第一次迭代考察低 2 位，分别把键分入低 2 位
为 00、01、10、11 的四个桶。第二次考察高 2 位，再根据高 2 位分入四个桶。与 1 位
示例一样，每个桶内都保留前一次迭代形成的键次序，因此第二次迭代后，所有键就按全部
4 位完全有序。

\begin{figure}[htbp]
  \centering
  \tiny
  \setlength{\tabcolsep}{1.5pt}
  \begin{tabular}{c|*{16}{c}}
    输入&1100&0011&0110&1001&1111&1000&0101&1010&1001&0110&1011&1101&0100&1010&0111&0000\\
    低 2 位&1100&1000&0100&0000&1001&0101&1001&1101&0110&1010&0110&1010&0011&1111&1011&0111\\
    高 2 位&0000&0011&0100&0101&0110&0110&0111&1000&1001&1001&1010&1010&1011&1100&1101&1111
  \end{tabular}
  \caption{使用 2 位基数的基数排序示例}
  \label{fig:14-10}
\end{figure}

与 1 位示例相似，可以为输入列表中的每个键分配一个线程，让线程求键的目标下标并写入
输出列表，从而并行化每次迭代。为优化内存合并，每个线程块可以先在共享内存中局部排序
自己的键，再以合并方式把局部桶写入全局内存。图 \ref{fig:14-11} 给出这种方式的
概念示例。

\begin{figure}[htbp]
  \centering
  \scriptsize
  \begin{tabular}{c}
    \fbox{每个线程块在共享内存中处理各自的 4 个键}\\[0.3em]
    $\downarrow$ 连续执行两轮局部 1 位基数排序 $\downarrow$\\[0.3em]
    \fbox{每块形成 00、01、10、11 四个稳定局部桶}\\[0.3em]
    $\downarrow$ 对局部桶大小表执行全局排他扫描 $\downarrow$\\[0.3em]
    \fbox{把各局部桶以连续区段写到四个全局桶}
  \end{tabular}
  \caption{用共享内存并行化并优化 2 位基数的一次基数排序迭代}
  \label{fig:14-11}
\end{figure}

1 位示例与 2 位示例的关键区别，是如何把键分入四个桶而非两个。在线程块内部局部
排序时，可以连续应用两次 1 位基数排序迭代来完成 2 位基数排序。两次 1 位迭代各自
需要一次排他扫描，但扫描都局限在线程块内，两次迭代之间无须跨线程块协调。一般而言，
对 $r$ 位基数，需要 $r$ 次局部 1 位迭代，把键分入 $2^r$ 个局部桶。

局部排序结束后，每个线程块必须确定自己的每个局部桶在全局输出列表中的位置。图
\ref{fig:14-12} 展示了 2 位示例的求法，其过程与图 \ref{fig:14-9} 的 1 位示例
相似。每块把各局部桶的键数存入表格，再扫描表格，取得各局部桶的全局位置。主要区别
是每块有四个局部桶而不是两个，因此排他扫描处理四行而不是两行。一般而言，$r$ 位基数
的排他扫描表格有 $2^r$ 行。

\begin{figure}[htbp]
  \centering
  \small
  \setlength{\tabcolsep}{7pt}
  \begin{tabular}{c|cccc}
    线程块&0&1&2&3\\\hline
    局部 00 桶大小&1&1&0&2\\
    局部 01 桶大小&1&1&2&0\\
    局部 10 桶大小&1&1&1&1\\
    局部 11 桶大小&1&1&1&1\\\hline
    \multicolumn{5}{c}{$\Downarrow$ 行主序展开后执行排他扫描 $\Downarrow$}\\\hline
    00 桶目标起点&0&1&2&2\\
    01 桶目标起点&4&5&6&8\\
    10 桶目标起点&8&9&10&11\\
    11 桶目标起点&12&13&14&15
  \end{tabular}
  \caption{确定 2 位基数下每个线程块局部桶的目标位置}
  \label{fig:14-12}
\end{figure}

使用更大基数的优点，是减少完全排好全部键所需的迭代次数。迭代减少意味着网格启动、
全局内存访问和全局排他扫描操作都更少。不过，更大基数也有缺点。第一，每个线程块拥有
更多局部桶，每个桶中的键更少。因此，各块需要写入更多不同的全局内存桶区段，而每个
区段的数据更少；基数越大，内存合并机会越少。第二，基数越大，应用全局排他扫描的表格
越大，扫描开销也越高。因此，基数不能任意增大。选择基数值时，必须在迭代次数与内存
合并行为、全局排他扫描开销之间取得平衡。多位基数排序的实现留作练习。

\section{用线程粗化改进内存合并}
\label{sec:radix-thread-coarsening}

在许多线程块间并行执行基数排序的代价之一，是全局内存写入的合并效果较差。每个线程块
都要把自己的局部桶写入全局内存。线程块越多，每块负责的键越少，局部桶也越小，写入
全局内存时暴露的合并机会便越少。如果这些线程块并行执行，这种代价或许值得；但如果
硬件本来就要把它们串行化，这笔代价就没有必要。

可以应用线程粗化解决这个问题：每个线程不再只负责输入列表中的一个键，而是负责多个键。
图 \ref{fig:14-13} 展示了如何对 2 位基数的一次迭代应用线程粗化。与图
\ref{fig:14-11} 相比，每个线程块负责更多键，因而每块的局部桶更大，暴露出更多内存
合并机会。比较两图可以看出，图 \ref{fig:14-13} 中连续线程更有可能写入连续内存位置。

\begin{figure}[htbp]
  \centering
  \scriptsize
  \begin{tabular}{c}
    \fbox{粗化块 0：处理输入前 8 个键}\qquad
    \fbox{粗化块 1：处理输入后 8 个键}\\[0.35em]
    $\downarrow$ 每个线程负责多个键，在共享内存中形成较大局部桶 $\downarrow$\\[0.35em]
    \fbox{块 0 的 00、01、10、11 桶}\qquad
    \fbox{块 1 的 00、01、10、11 桶}\\[0.35em]
    $\downarrow$ 更少、更大的连续区段写入全局内存 $\downarrow$\\[0.35em]
    \fbox{按低 2 位稳定划分后的全局输出列表}
  \end{tabular}
  \caption{对 2 位基数应用线程粗化，以改进内存合并}
  \label{fig:14-13}
\end{figure}

在许多线程块间并行化基数排序的另一项开销，是为了确定各线程块局部桶的目标位置而执行
全局排他扫描。回忆图 \ref{fig:14-12}，扫描表格的大小与桶数和线程块数都成正比。
应用线程粗化会减少线程块数，从而缩小表格，降低排他扫描开销。在线程粗化的基数排序
实现留作练习。

\section{其他并行排序方法}
\label{sec:other-parallel-sorts}

本章概述的算法只是并行排序数据的许多方法中的一部分。本节简要介绍读者可能感兴趣的
其他方法。

第 \ref{sec:parallel-odd-even-sort} 节的奇偶换位排序采用固定比较模式，在元素顺序
错误时交换它们。它很容易并行化，因为每一步比较的键对互不重叠。有一整类排序方法都
使用固定比较模式对序列排序，而且常常并行执行；它们通常称为\textbf{排序网络
（sorting network）}。最著名的并行排序网络是 Batcher 的双调排序与奇偶归并排序
\cite{batcher1968networks}。Batcher 算法处理固定长度序列，比奇偶换位排序更高效；
对 $N$ 个元素只需 $O(N\log^2N)$ 次比较。尽管这一渐近代价劣于归并排序等算法的
$O(N\log N)$，实践中它们凭借简单性，经常成为短序列上最高效的方法。

大多数不采用排序网络固定比较集合的并行比较排序，可以分为两大类。第一类把未排序
输入划分为数据片，分别排好每片，再把大部分工作放在组合这些数据片以形成输出上。本章
的归并排序就是例子：大部分工作发生于组合各有序数据片的归并树中。第二类把大部分工作
放在划分未排序序列上，使最后组合各分区相对容易。\textbf{采样排序（sample sort）}
\cite{frazer1970samplesort}是这一类的典型例子。采样排序先从输入中选择 $p-1$ 个键
（例如随机选择）并排好它们，再用这些键把输入划分成 $p$ 个桶，使桶 $k$ 中的所有键
都大于任意桶 $j<k$ 中的键，并小于任意桶 $j>k$ 中的键。这一步可以看作快速排序二路
划分向 $p$ 路的推广。如此划分数据后，各桶可以独立排序，最后只需按顺序连接各桶便能
形成有序输出。

对于极大的序列，尤其数据必须分布在多个物理内存中，甚至分布到同一节点中多块 GPU
的内存时，采样排序通常是最高效的选择。实践中经常对键进行过采样，因为适度过采样能
以很高概率得到均衡分区\cite{blelloch1991comparison}。

正如归并排序和采样排序分别代表比较排序中的自底向上与自顶向下策略，基数排序也能按
这两类策略设计。本章介绍的基数排序更完整地说是\textbf{最低有效位（LSB）基数排序}，
更一般地称为\textbf{最低有效数位（LSD）基数排序}。算法从键的最低有效数位开始，
逐步处理到最高有效数位。\textbf{最高有效数位（MSD）基数排序}采用相反策略：先用
最高有效数位把输入划分成对应于该位各可能值的桶，再在每个桶内独立使用下一个最高
有效数位继续划分；处理到最低有效数位时，整个序列便已排好。

MSD 基数排序与采样排序一样，通常更适合超大序列。LSD 基数排序每一步都需要全局重排
数据，而 MSD 基数排序的每一步会在逐渐变得更局部的数据区域上运行。

\section{小结}
\label{sec:sorting-summary}

本章介绍了如何在 GPU 上并行排序键及其关联值。首先介绍两种基于比较的并行排序算法：
并行奇偶排序和并行归并排序。并行奇偶排序在奇数对与偶数对之间交替，反复交换列表中
次序错误的相邻元素对，直到列表有序，从而避免竞争条件。并行归并排序并行执行不同输入
区段间的独立归并操作，同时利用第 13 章的并行归并方法，在每个归并操作内部继续并行化。

随后，本章重点讨论工作复杂度低于 $O(N\log N)$ 的非比较算法基数排序。基数排序反复
把键分配到各桶中，对键的每个数位都重复分配，并保留上次数位迭代形成的次序，最终保证
键按全部数位有序。每次迭代为输入列表中的每个键分配一个线程，让线程与其他线程协作
执行排他扫描，找出该键在输出列表中的目标位置。

优化基数排序的一项关键挑战，是写入输出列表时实现合并内存访问。一项重要优化，是让
每个线程块先在共享内存的局部桶中执行局部排序，再以合并方式把每个局部桶写入全局内存。
另一项优化是增大基数，减少所需迭代数和网格启动数；但基数不能过大，否则会恶化内存
合并，并增加全局排他扫描开销。最后，线程粗化既能改进内存合并，也能降低全局排他扫描
开销。

在 GPU 上实现并优化并行排序算法很复杂，建议程序员使用 Thrust
\cite{bell2012thrust-sort,nvidia-thrust-sort} 或 CUB\cite{nvidia-cub-sort} 等 GPU
并行排序库，而不是从头编写排序内核。即便如此，并行排序仍是研究并行模式优化权衡的
一个有趣案例。

\phantomsection
\section*{练习}
\addcontentsline{toc}{section}{练习}

\begin{enumerate}[label=\arabic*.,leftmargin=2.7em]
  \item 扩展图 \ref{fig:14-7} 的内核，使用共享内存改进内存合并。
  \item 扩展图 \ref{fig:14-7} 的内核，使其支持多位基数。
  \item 扩展图 \ref{fig:14-7} 的内核，应用线程粗化改进内存合并。
  \item 使用第 13 章的并行归并实现并行归并排序。
\end{enumerate}

\begin{chapterbibliography}{9}
  \bibitem{satish2009sorting}
  N. Satish, M. Harris, M. Garland, Designing efficient sorting algorithms for
  manycore GPUs, in: \emph{2009 IEEE International Symposium on Parallel \&
  Distributed Processing}, IEEE, 2009, pp. 1--10.

  \bibitem{cole1988merge}
  R. Cole, Parallel merge sort, \emph{SIAM Journal on Computing} 17 (4) (1988)
  770--785.

  \bibitem{adinets2022onesweep}
  A. Adinets, D. Merrill, OneSweep: a faster least significant digit radix sort
  for GPUs, arXiv preprint arXiv:2206.01784, 2022.

  \bibitem{batcher1968networks}
  K.E. Batcher, Sorting networks and their applications, in: \emph{Proceedings
  of the April 30--May 2, 1968, Spring Joint Computer Conference}, 1968,
  pp. 307--314.

  \bibitem{frazer1970samplesort}
  W.D. Frazer, A.C. McKellar, Samplesort: a sampling approach to minimal storage
  tree sorting, \emph{Journal of the ACM} 17 (3) (1970) 496--507.

  \bibitem{blelloch1991comparison}
  G.E. Blelloch, C.E. Leiserson, B.M. Maggs, C.G. Plaxton, S.J. Smith,
  M. Zagha, A comparison of sorting algorithms for the Connection Machine CM-2,
  in: \emph{Proceedings of the Third Annual ACM Symposium on Parallel Algorithms
  and Architectures}, 1991, pp. 3--16.

  \bibitem{bell2012thrust-sort}
  N. Bell, J. Hoberock, Thrust: a productivity-oriented library for CUDA,
  in: \emph{GPU Computing Gems Jade Edition}, Elsevier, 2012, pp. 359--371.

  \bibitem{nvidia-thrust-sort}
  NVIDIA, Thrust: the C++ parallel algorithms library, 2025.
  \url{https://nvidia.github.io/cccl/thrust/index.html}.

  \bibitem{nvidia-cub-sort}
  NVIDIA, CUB, 2025.
  \url{https://nvidia.github.io/cccl/cub/index.html}.
\end{chapterbibliography}
